\documentclass[a4paper,fontset=none]{ctexart}

\usepackage{indentfirst}
\setlength{\parindent}{0em}
\usepackage{autobreak}
\allowdisplaybreaks
\usepackage{parskip}
\usepackage{geometry}
\geometry{top=3cm,bottom=3cm,left=2cm,right=2cm}
\usepackage{anyfontsize}

\usepackage{fontspec}
\setmainfont{STIX Two Text}
\usepackage{unicode-math}
\unimathsetup{mathrm=sym,mathbf=sym,mathsf=sym,bold-style=ISO}
\setmathfont{STIX Two Math}
\usepackage{physics2}
\usephysicsmodule{ab,op.legacy}
\usepackage{fixdif,derivative}
\newcommand{\um}{\upmu\mathrm{m}}
\newcommand{\ang}{\text{\r{A}}}
\AtBeginDocument{
    \let\uglyepsilon\epsilon
    \renewcommand\epsilon{\varepsilon}
}

\setCJKmainfont{Source Han Serif CN}[BoldFont=Source Han Sans CN Medium,ItalicFont=FZKai-Z03]

\title{\textbf{星际介质天文学作业}}
\author{张植竣}
\date{2024年5月3日}

\begin{document}

\maketitle

\begin{enumerate}
    \item[32.1 (a)]
    巨分子云的总质量为：
    \begin{align*}
        M_{\mathrm{tot}}
        &= \int_{10^3\,M_\odot}^{M_u} M\d{N} \\
        &= \int_{10^3\,M_\odot}^{M_u} M\odv{N}{M} \d{M} \\
        &= \int_{10^3\,M_\odot}^{M_u} \odv{N}{\ln{M}} \d{M} \\
        &= \int_{10^3\,M_\odot}^{M_u} N_u \ab(\frac{M}{M_u})^{-\alpha} \d{M} \\
        &= \int_{10^3}^{6 \times 10^6} 63 \times \ab(\frac{x}{6 \times 10^6})^{-0.6}\,M_\odot \d{x},\quad x = \frac{M}{M_\odot} \\
        &= 9.16 \times 10^8\,M_\odot
    \end{align*}

    \medskip
    \item[32.1 (b)]
    质量大于$10^6\,M_\odot$的巨分子云的数量为：
    \begin{align*}
        N(M > 10^6\,M_\odot)
        &= \int_{10^6\,M_\odot}^{M_u} \d{N} \\
        &= \int_{10^6\,M_\odot}^{M_u} \odv{N}{M} \d{M} \\
        &= \int_{10^6\,M_\odot}^{M_u} \odv{N}{\ln{M}} \frac{\d{M}}{M} \\
        &= \int_{10^6\,M_\odot}^{M_u} \frac{N_u}{M} \ab(\frac{M}{M_u})^{-\alpha} \d{M} \\
        &= \int_{10^6}^{6 \times 10^6} \frac{63}{x} \times \ab(\frac{x}{6 \times 10^6})^{-0.6} \d{x},\quad x = \frac{M}{M_\odot} \\
        &= 203
    \end{align*}

\end{enumerate}

\end{document}